% ============================================================================
% Indonesian UI and structural strings for the One Chemistry Book series.
% Loaded by styles/onechemistry.sty when \booklang is id.
%
% Standard Bahasa Indonesia, EYD spelling. Indonesian marks no plural on the
% noun, so every "plural" name below is deliberately identical to its singular
% — that is the correct form, not an oversight.
% ============================================================================

% Theorem / statement environment titles
\newcommand{\omnameDefinition}{Definisi}
\newcommand{\omnameTheorem}{Teorema}
\newcommand{\omnameProposition}{Proposisi}
\newcommand{\omnameLemma}{Lema}
\newcommand{\omnameCorollary}{Akibat}
\newcommand{\omnameMethod}{Metode}
\newcommand{\omnameExample}{Contoh}
\newcommand{\omnameNotation}{Notasi}
\newcommand{\omnameRemark}{Catatan}
\newcommand{\omnameExercise}{Latihan}
\newcommand{\omnameExercises}{Latihan}
\newcommand{\omnameProblem}{Soal}
\newcommand{\omnameProblems}{Soal}
\newcommand{\omnameProof}{Bukti}

% Admitted results and solutions appendix headers
\newcommand{\omadmittedtext}{Diterima tanpa bukti pada tingkat ini.}
% Used as "\omsolutionof~\cref{exo:...}": Indonesian needs no preposition here.
\newcommand{\omsolutionof}{Solusi}
\newcommand{\ompage}{halaman}
\newcommand{\omnameGotoSolution}{Solusi hlm.}

% Misc text macros
\newcommand{\st}{\text{ sedemikian sehingga }}

% Cover / title page
\newcommand{\bookmaintitle}{One Chemistry Book}
\newcommand{\booktagline}{Satu Buku Kimia untuk Menguasai Semuanya}
\newcommand{\bookdescription}{Kimia dari kelas satu sekolah dasar hingga akhir program sarjana}
\newcommand{\bookcontributors}{Para kontributor One Chemistry Book}
\newcommand{\bookversionlabel}{Versi}

% Colophon
\newcommand{\omcolophonsource}{Kode sumber, laporan kesalahan dan kontribusi:}
\newcommand{\omcolophonlicense}{Kecuali dinyatakan lain, buku ini dilisensikan di bawah Creative Commons CC BY-NC-SA 4.0:}
\newcommand{\omcolophoncredit}{Cantumkan sumber sebagai: One Chemistry Book, One Course (one-course.com).}

% Chemistry boxes (styles/onechemistry.sty): the recall box that opens a
% chapter building on another one, the lab box, the history box, the safety
% box. Placeholder wording by the coordinator (2026-10-04); the Book 1
% `id` agent owns it and settles the register (the recall title addresses
% the reader in the same register as the book's exercise stems).
\newcommand{\omnameRecall}{Yang sudah kamu ketahui}
\newcommand{\omnameInTheLab}{Di laboratorium}
\newcommand{\omnameHistory}{Sejarah}
\newcommand{\omnameSafety}{Keselamatan}

% Solutions appendix part title
\newcommand{\omsolutionspart}{Solusi Latihan}

% Part titles (Kelas 1–12; parallel to English Grade N)
\@namedef{omstr@part.grade1}{Kelas 1 (Usia 6--7 tahun)}
\@namedef{omstr@part.grade2}{Kelas 2 (Usia 7--8 tahun)}
\@namedef{omstr@part.grade3}{Kelas 3 (Usia 8--9 tahun)}
\@namedef{omstr@part.grade4}{Kelas 4 (Usia 9--10 tahun)}
\@namedef{omstr@part.grade5}{Kelas 5 (Usia 10--11 tahun)}
\@namedef{omstr@part.grade6}{Kelas 6 (Usia 11--12 tahun)}
\@namedef{omstr@part.grade7}{Kelas 7 (Usia 12--13 tahun)}
\@namedef{omstr@part.grade8}{Kelas 8 (Usia 13--14 tahun)}
\@namedef{omstr@part.grade9}{Kelas 9 (Usia 14--15 tahun)}
\@namedef{omstr@part.grade10}{Kelas 10 (Usia 15--16 tahun)}
\@namedef{omstr@part.grade11}{Kelas 11 (Usia 16--17 tahun)}
\@namedef{omstr@part.grade12}{Kelas 12 (Usia 17--18 tahun)}

% Part titles (Sarjana Tahun 1–3; parallel to English Bachelor Year N).
\@namedef{omstr@part.bachelor1}{Sarjana Tahun ke-1 (Usia 18--19 tahun)}
\@namedef{omstr@part.bachelor2}{Sarjana Tahun ke-2 (Usia 19--20 tahun)}
\@namedef{omstr@part.bachelor3}{Sarjana Tahun ke-3 (Usia 20--21 tahun)}

% Plural names + \omnameChapter, used by the \crefname declarations in
% onechemistry.sty. Identical to the singulars on purpose (see the header).
\newcommand{\omnameDefinitions}{Definisi}
\newcommand{\omnameTheorems}{Teorema}
\newcommand{\omnamePropositions}{Proposisi}
\newcommand{\omnameLemmas}{Lema}
\newcommand{\omnameCorollarys}{Akibat}
\newcommand{\omnameMethods}{Metode}
\newcommand{\omnameExamples}{Contoh}
\newcommand{\omnameNotations}{Notasi}
\newcommand{\omnameRemarks}{Catatan}
\newcommand{\omnameChapter}{Bab}
\newcommand{\omnameChapters}{Bab}
% Section, for \cref{sec:…} (2026-10-04: it had no \crefname, so cleveref
% printed its English built-in "Section" in every edition).
\newcommand{\omnameSection}{Subbab}
\newcommand{\omnameSections}{Subbab}

% siunitx and cleveref keep their English conjunctions unless told otherwise:
% \qtyrange printed "5 to 4186 Hz" and \cref{a,b} printed "Bab 25 and 29" in
% the middle of Indonesian prose.
% cleveref installs its conjunctions at \begin{document}, so redefine them there.
% siunitx typesets range phrases and list separators in MATH MODE, where a
% bare string loses its spaces (\qtyrange inside $...$ printed "1a100") and a
% non-ASCII character is worse: "a" expands to \`a, which LaTeX reports as
% "Command \` invalid in math mode" and drops from the page. siunitx's own
% defaults are wrapped in \text{} for exactly this reason; these overrides
% must be too. Found on Biology Book 3 fr, 2026-09-05, where the book's 56
% \qtyrange sites made it visible; it was latent in every language edition of
% every book in this repo.
\sisetup{range-phrase={\text{ sampai }}, list-pair-separator={\text{ dan }},
         list-final-separator={\text{ dan }}}
\AtBeginDocument{%
  \renewcommand{\crefpairconjunction}{ dan }%
  \renewcommand{\creflastconjunction}{ dan }%
  \renewcommand{\crefrangeconjunction}{ sampai }%
  % a \cref spanning two label types builds a group list, which has its
  % own pair of conjunctions: "Teorema 2.14 dan 2.18 and Lema 2.17".
  \renewcommand{\crefpairgroupconjunction}{ dan }%
  \renewcommand{\creflastgroupconjunction}{ dan }%
}

% \today is English without babel (which is optional here), so the cover date
% read "August 20, 2026" on an Indonesian book. Indonesian writes it as
% "20 Agustus 2026", month capitalised.
\renewcommand{\today}{\number\day~\ifcase\month\or Januari\or Februari\or
  Maret\or April\or Mei\or Juni\or Juli\or Agustus\or September\or Oktober\or
  November\or Desember\fi\space\number\year}

% LaTeX's own structural names. babel-indonesian sets the same words, but it
% is optional here (see onechemistry.sty), so declare them anyway.
\renewcommand{\chaptername}{Bab}
\renewcommand{\partname}{Bagian}
\renewcommand{\contentsname}{Daftar Isi}
\renewcommand{\appendixname}{Lampiran}
\renewcommand{\indexname}{Indeks}

% Periodic-table legend (\omperiodictable[legend], styles/onechemistry.sty).
% Coordinator wording, 2026-10-04 (it used to be hard-coded English in the
% style file, so no edition could translate it); the Book 1 agent of this
% language checks it against the book's own terminology.
\newcommand{\omnamePTmetal}{logam}
\newcommand{\omnamePTmetalloid}{metaloid}
\newcommand{\omnamePTnonmetal}{nonlogam}
\newcommand{\omnamePTs}{blok s}
\newcommand{\omnamePTp}{blok p}
\newcommand{\omnamePTd}{blok d}
\newcommand{\omnamePTf}{blok f}
\newcommand{\omnamePTalkali}{logam alkali}
\newcommand{\omnamePTalkaline}{logam alkali tanah}
\newcommand{\omnamePThalogen}{halogen}
\newcommand{\omnamePTnoble}{gas mulia}
